% ====================================================================
% Capítulo 11: Métodos numéricos
% ====================================================================
\chapter{Métodos numéricos}
\label{ch:metodos_numericos}

Este capítulo describe en detalle los algoritmos numéricos utilizados en la 
suite \texttt{Mpemba-X v2}. Cada método se documenta con su análisis de 
convergencia, su orden de exactitud, las consideraciones de estabilidad y su 
implementación paralela.

\section{Integración temporal de ecuaciones diferenciales ordinarias: Runge-Kutta de 4º orden}

\subsection{Formulación}

Para una ecuación diferencial $\dot{\vb{y}} = \vb{f}(t, \vb{y})$, el integrador 
Runge-Kutta de cuarto orden (RK4) construye la actualización 
$\vb{y}(t + h) = \vb{y}(t) + h \vb{\Phi}(t, \vb{y}, h)$ con:
\begin{equation}
\begin{aligned}
    \vb{k}_1 &= \vb{f}(t, \vb{y}) \\
    \vb{k}_2 &= \vb{f}(t + h/2, \vb{y} + h \vb{k}_1 / 2) \\
    \vb{k}_3 &= \vb{f}(t + h/2, \vb{y} + h \vb{k}_2 / 2) \\
    \vb{k}_4 &= \vb{f}(t + h, \vb{y} + h \vb{k}_3) \\
    \vb{\Phi} &= \tfrac{1}{6}(\vb{k}_1 + 2 \vb{k}_2 + 2 \vb{k}_3 + \vb{k}_4)
\end{aligned}
\end{equation}

\subsection{Error de truncamiento}

El error global del paso es $\mathcal{O}(h^4)$. Para sistemas suaves, esta 
exactitud es suficiente con $h \sim 10^{-3}$. RK4 es \emph{condicionalmente 
estable}: para sistemas lineales $\dot{\vb{y}} = A \vb{y}$, la región de 
estabilidad es el disco $|1 + h\lambda + (h\lambda)^2/2 + (h\lambda)^3/6 + (h\lambda)^4/24| \leq 1$ 
en el plano complejo de $h \lambda_{\mathrm{eig}(A)}$.

\subsection{Aplicación a la ecuación maestra markoviana}

En el módulo \texttt{mpemba\_markovian}, RK4 se aplica a la ecuación
$\dot{\vb{p}} = \vb{R} \vb{p}$, donde $\vb{R}$ es la matriz de tasas del 
\cref{eq:master}. Cada producto $\vb{R}\vb{p}$ se realiza con OpenMP en el lazo 
sobre las componentes de $\vb{p}$:
\begin{lstlisting}[language=C++,caption={Multiplicación matriz-vector paralelizada}]
void dense_matvec(const DenseMarkov& M, const double* x, double* y) {
    int n = M.n;
    #pragma omp parallel for schedule(static)
    for (int i = 0; i < n; ++i) {
        double s = 0.0;
        const double* row = &M.R[(std::size_t)i * n];
        for (int j = 0; j < n; ++j) s += row[j] * x[j];
        y[i] = s;
    }
}
\end{lstlisting}

\section{Integración estocástica: Euler-Maruyama y esquemas relacionados}

\subsection{Esquema Euler-Maruyama}

Para la ecuación de Langevin $\dd x = a(x) \dd t + b(x) \dd W_t$ donde $W_t$ es 
proceso de Wiener, el esquema Euler-Maruyama da
\begin{equation}
    x(t + h) = x(t) + a(x(t)) h + b(x(t)) \sqrt{h} \, \eta_t, 
    \quad \eta_t \sim \mathcal{N}(0, 1).
    \label{eq:em_method}
\end{equation}
%
Es un esquema de orden \emph{débil} 1 (en la convergencia de momentos) y orden 
\emph{fuerte} $1/2$ (en la convergencia de trayectorias).

\subsection{Condiciones de estabilidad}

Para potenciales armónicos $U = \tfrac{1}{2}\omega^2 x^2$ con dinámica de 
Langevin sobreamortiguada, el esquema EM es estable si
\begin{equation}
    h < \frac{2}{\omega^2}.
\end{equation}
%
\noindent En el módulo \texttt{mpemba\_langevin}, $h = 5 \times 10^{-5}$ con 
$U_0 = 6$, $h = 0.8$ asegura estabilidad y convergencia adecuada.

\subsection{Generación de números aleatorios}

La generación de variables normales requiere un PRNG de alta calidad y 
\emph{thread-safe}. \texttt{Mpemba-X v2} implementa Philox-4x32-10 de 
Salmon \emph{et al.}~\cite{salmon2011}:

\begin{lstlisting}[language=C++,caption={Esqueleto del PRNG Philox}]
struct PhiloxRNG {
    uint64_t key, counter;
    uint32_t buffer[4];
    int buffer_pos = 4;
    
    void refill() {
        // 10 rondas de la función de mezcla de Philox
        // ... implementación detallada en core/philox_rng.hpp
        counter++;
    }
    
    double uniform() {
        if (buffer_pos == 4) { refill(); buffer_pos = 0; }
        return ldexp(buffer[buffer_pos++], -32);
    }
    
    double normal() {
        // Box-Muller
        double u1 = uniform(), u2 = uniform();
        return std::sqrt(-2.0 * std::log(u1)) * std::cos(2 * M_PI * u2);
    }
};
\end{lstlisting}

\noindent Cada hilo OpenMP recibe un PRNG con \texttt{counter} 
deterministicamente diferente, garantizando reproducibilidad y ausencia de 
correlación entre hilos.

\section{Discretización espacial: diferencias finitas y volúmenes finitos}

\subsection{Ecuación de difusión 1D con esquema explícito}

El módulo \texttt{mpemba\_water\_fourier} resuelve la 
\cref{eq:fourier_zhang}. Se discretiza el dominio en $N$ celdas de tamaño 
$\Delta x$. Para la celda $i$ interior, la actualización es:
\begin{equation}
\begin{aligned}
    q^{n}_{i-1/2} &= -\kappa_{i-1/2} \frac{\theta^{n}_{i} - \theta^{n}_{i-1}}{\Delta x}, \\
    q^{n}_{i+1/2} &= -\kappa_{i+1/2} \frac{\theta^{n}_{i+1} - \theta^{n}_{i}}{\Delta x}, \\
    \theta^{n+1}_{i} &= \theta^{n}_{i} + \frac{\Delta t}{\rho_{i} c_{p,i} \Delta x} 
                       \left[q^{n}_{i-1/2} - q^{n}_{i+1/2}\right] + \cdots
\end{aligned}
\end{equation}
%
donde $\kappa_{i\pm1/2}$ es la \emph{media armónica} entre las celdas vecinas:
\begin{equation}
    \kappa_{i+1/2} = \frac{2 \kappa_{i} \kappa_{i+1}}{\kappa_{i} + \kappa_{i+1}}.
\end{equation}

\subsection{Condición CFL}

El esquema explícito es condicionalmente estable. Para difusión pura, el límite 
CFL es:
\begin{equation}
    \Delta t \leq \frac{\Delta x^2}{2 \alpha_{\max}},
\end{equation}
%
con $\alpha = \kappa/(\rho c_p)$ la difusividad térmica. En presencia de 
condición de frontera Robin con coeficiente $h$, la restricción más estricta 
proviene de la celda de borde:
\begin{equation}
    \Delta t \leq \frac{0.5 \cdot \rho_0 c_{p,0} \Delta x^2}{\kappa_0 + h \Delta x}.
    \label{eq:cfl_robin}
\end{equation}
%
\noindent El módulo \texttt{mpemba\_water\_fourier} usa un \emph{sub-stepping 
adaptativo}: dado un paso externo $\Delta t_{\mathrm{out}}$ (por ejemplo \SI{5}{\second} 
para registro de datos), se calcula el límite CFL local $\Delta t_{\mathrm{CFL}}$ 
y se realizan $\lceil \Delta t_{\mathrm{out}} / \Delta t_{\mathrm{CFL}} \rceil$ 
sub-pasos.

\section{Diagonalización de matrices simétricas: método de Jacobi}

Para el módulo \texttt{mpemba\_klich\_raz} se requiere la diagonalización 
completa de la matriz simétrica $\tilde{\vb{R}}$ del \cref{eq:R_tilde}. Se 
implementa el algoritmo clásico de \emph{rotaciones de Jacobi}~\cite{golub2013}.

\subsection{Algoritmo}

Dada una matriz simétrica $A \in \mathbb{R}^{n \times n}$:
\begin{enumerate}
\item Encontrar el par $(p, q)$ con $\abs{A_{pq}}$ máximo fuera de la diagonal.
\item Calcular el ángulo de rotación $\theta$ tal que la transformación de 
Givens $G(p, q, \theta)$ anule $A_{pq}$:
\begin{equation}
    \tan(2\theta) = \frac{2 A_{pq}}{A_{qq} - A_{pp}}.
\end{equation}
\item Actualizar $A \leftarrow G^{\top} A G$ y el acumulador de vectores 
propios $V \leftarrow V G$.
\item Repetir hasta que la norma de Frobenius de la parte fuera de la diagonal 
sea menor que una tolerancia $\epsilon$.
\end{enumerate}

\subsection{Complejidad y convergencia}

El método tiene complejidad $\mathcal{O}(n^3)$ por barrido completo y converge 
cuadráticamente cerca de la convergencia. Para los tamaños usados en 
Klich-Raz ($n = L = 10$ a $14$), el costo total es del orden de $10^{-3}$ 
segundos por realización.

\section{Direct Simulation Monte Carlo (DSMC) para gases granulares}

El módulo \texttt{mpemba\_granular\_dsmc} implementa el método DSMC clásico de 
Bird~\cite{bird1994} para simular un gas granular inelástico homogéneo.

\subsection{Algoritmo NTC (No-Time Counter)}

A cada paso de tiempo $\Delta t$:
\begin{enumerate}
\item Avance balístico: $\vb{r}_i \to \vb{r}_i + \vb{v}_i \Delta t$ (en sistemas 
homogéneos sin gradientes, este paso es trivial).
\item Selección de pares de colisión: número de pares candidatos
\begin{equation}
    N_{\mathrm{pairs}} = \frac{N (N-1)}{2} \sigma_T \langle \abs{\vb{v}_{\mathrm{rel}}} \rangle_{\max} \frac{\Delta t}{V},
\end{equation}
donde $\sigma_T = \pi (2a)^2$ es la sección eficaz total y $V$ el volumen.
\item Para cada par candidato $(i, j)$: aceptar con probabilidad
\begin{equation}
    P_{\mathrm{accept}} = \frac{\abs{\vb{v}_i - \vb{v}_j}}{\langle \abs{\vb{v}_{\mathrm{rel}}} \rangle_{\max}}.
\end{equation}
\item Para pares aceptados: aplicar la regla inelástica de \cref{eq:inelastic_collision}.
\end{enumerate}

\subsection{Paralelización OpenMP de DSMC}

La paralelización es no trivial porque las colisiones potencialmente involucran 
pares globales. La estrategia de \texttt{Mpemba-X v2} divide el sistema en 
\emph{sub-celdas} espaciales (aunque homogéneo, esto es una conveniencia 
algorítmica) y procesa pares dentro de cada sub-celda en paralelo. Para 
sistemas pequeños ($N \leq 10^4$), se usa la implementación serial.

\section{Dinámica molecular: Velocity-Verlet y termostato de Berendsen}

\subsection{Velocity-Verlet}

Para un sistema de $N$ partículas con masas $m_i$ y fuerzas $\vb{F}_i$, el 
integrador Velocity-Verlet es:
\begin{equation}
\begin{aligned}
    \vb{v}_i(t + h/2) &= \vb{v}_i(t) + \frac{h}{2} \vb{F}_i(t)/m_i, \\
    \vb{r}_i(t + h) &= \vb{r}_i(t) + h \, \vb{v}_i(t + h/2), \\
    \vb{F}_i(t + h) &= \text{recalcular en } \vb{r}_i(t+h), \\
    \vb{v}_i(t + h) &= \vb{v}_i(t + h/2) + \frac{h}{2} \vb{F}_i(t+h)/m_i.
\end{aligned}
\end{equation}
%
El esquema es simpléctico, conserva energía hasta orden $\mathcal{O}(h^2)$ y 
es estable para $h < 2/\omega_{\max}$ donde $\omega_{\max}$ es la frecuencia 
más alta del sistema.

\subsection{Termostato de Berendsen}

Para mantener un ensemble $NVT$ se acopla un \emph{termostato de 
Berendsen}~\cite{berendsen1984} que reescala las velocidades:
\begin{equation}
    \vb{v}_i \to \vb{v}_i \cdot \lambda(T_{\mathrm{cur}}), 
    \qquad \lambda = \sqrt{1 + \frac{h}{\tau_T}\left(\frac{T_{\mathrm{target}}}{T_{\mathrm{cur}}} - 1\right)},
\end{equation}
%
donde $\tau_T$ es la constante temporal del termostato (típicamente 
$\tau_T \sim 100 h$). Este termostato no genera el ensemble $NVT$ canónico 
exactamente (no satisface el balance detallado), pero es robusto y simple, 
adecuado para simulaciones de \emph{quench}.

\subsection{Potencial mW para agua}

El módulo \texttt{mpemba\_water\_md} usa el potencial mW (\emph{monatomic 
Water}) de Molinero \& Moore~\cite{molinero2009}, un Stillinger-Weber tetraédrico:
\begin{equation}
    U = \sum_{i,j} \phi_2(r_{ij}) + \sum_{i,j,k} \phi_3(r_{ij}, r_{ik}, \theta_{ijk}),
\end{equation}
%
con
\begin{align}
    \phi_2(r) &= A \epsilon \left[B \left(\frac{\sigma}{r}\right)^4 - 1\right] \exp\!\left[\frac{\sigma}{r - a \sigma}\right], \\
    \phi_3(r_1, r_2, \theta) &= \lambda \epsilon \left[\cos\theta - \cos\theta_0\right]^2 
    \exp\!\left[\frac{\gamma \sigma}{r_1 - a \sigma}\right] \exp\!\left[\frac{\gamma \sigma}{r_2 - a \sigma}\right],
\end{align}
%
con parámetros estándar $\sigma = \SI{2.3925}{\angstrom}$, 
$\epsilon = \SI{6.189}{\kilo\calorie\per\mole}$, $A = 7.0496$, $B = 0.6022$, 
$\gamma = 1.2$, $\lambda = 23.15$, $\theta_0 = 109.47°$.

\subsection{Implementación}

La implementación actual de \texttt{Mpemba-X v2} usa solo la parte 2-body de 
mW (la parte 3-body se omite para reducir el costo computacional al precio 
de una pérdida de fidelidad en la estructura tetraédrica). Para producción 
quantitativa se recomienda LAMMPS con mW completo.

\section{Paralelización híbrida OpenMP + MPI}

\subsection{Distribución MPI}

\texttt{Mpemba-X v2} usa MPI para distribuir trabajo entre nodos o procesos 
independientes. Patrones típicos:

\begin{itemize}
\item \textbf{Distribución de temperaturas iniciales}: cada rank ejecuta los 
quenches de un subconjunto de los $\Tinit$ del barrido. No requiere 
comunicación durante la simulación.
\item \textbf{Distribución de trayectorias estocásticas}: cada rank simula 
$N_{\mathrm{traj}}/k$ trayectorias independientes. La agregación final del 
histograma se realiza con \texttt{MPI\_Allreduce}.
\item \textbf{Distribución de realizaciones del desorden}: para Klich-Raz, 
cada rank procesa $N_{\mathrm{realiz}}/k$ realizaciones del modelo REM.
\end{itemize}

\subsection{Paralelización OpenMP intra-rank}

Dentro de cada rank MPI, OpenMP paraleliza los lazos críticos:
\begin{itemize}
\item Producto matriz-vector $\vb{R} \vb{p}$ en la ecuación maestra.
\item Lazo sobre trayectorias Langevin.
\item Lazo sobre pares de fuerzas en MD.
\end{itemize}
%
La cantidad de hilos por rank se configura mediante 
\texttt{OMP\_NUM\_THREADS}. Configuraciones típicas:
\begin{itemize}
\item Laptop (4 cores): 2 ranks $\times$ 2 hilos.
\item Workstation (16 cores): 4 ranks $\times$ 4 hilos.
\item Cluster (nodos de 32 cores): 8 ranks $\times$ 4 hilos por nodo.
\end{itemize}
